% Copyright (c) 2026 by C. Klukas.
% Licensed under the MIT License. See LICENSE for details.

A \paglet{} of the previous chapters lives on one host at a time; when that
host goes down, the \paglet{} is down with it (Section~\ref{sec:movement}).
A \emph{zombie \paglet{}} cannot die while one host of its mesh is alive,
which is what the name says: it keeps coming back. It has a replica on
every host. One replica, the \emph{main}, works; the
others sleep, hold the main's latest checkpoint and take over when the main
disappears. After a partition, the replicas agree on the one that got
furthest, and what both sides observed is merged. After the whole mesh was
down, the most advanced replica goes on. When the work is done, exactly one
result reaches the sink, and every replica ends. This chapter describes the
protocol, the SDK class that runs it, and what was checked about it.

\subsection{Ground rules}\label{sec:zombies-rules}

A replica that hears nothing cannot tell dead hosts from a cut network, so
the protocol does not try to keep a single main at all times. Four rules
follow~\cite{pagletsCppDocs}:

\begin{itemize}
\item \textbf{Several mains may run for a while.} A cut-off replica takes
  over after a timeout; during a partition each side may have its own main.
\item \textbf{One result counts.} Uniqueness is decided where it can be
  decided atomically: at the sink, which accepts the first result for a
  group and answers every later one with \emph{already delivered}.
\item \textbf{Mains that meet agree on one.} The one further along wins.
  The other forgets its work, sleeps and takes checkpoints from the winner.
  Its observations are not lost: journals are merged.
\item \textbf{Work repeats from the last checkpoint.} A replica that takes
  over continues from the checkpoint it holds, so work since then is done
  again (at least once). Side effects other than the result should be
  idempotent.
\end{itemize}

The protocol is a pure C++ state machine: the replica state, the messages,
and \code{step()}, the transition function of one replica. It makes no ABI
call and reads no clock. The same code is run by the guest SDK and explored
by the model checker and the simulator; recorded traces are checked against
its transition table and its order on ranks. The excerpts below are taken from its header,
\path{cpp/common/include/paglets/protocols/zombie/zombie.hpp}.

\subsection{Replicas, roles and ranks}\label{sec:zombies-roles}

A \emph{group} is one zombie \paglet{}: a 128-bit group ID, a sink, and a
replica on every host. A replica is named by its host, by the first 8 bytes
of the host's key ID (\code{replica\_id()}), so names are stable, the same
on every replica, and do not depend on which main made the replica. (Two
hosts whose key IDs share their first 16 hex digits would get the same
name; the SDK then makes no replica on the second one and logs a warning.)
A replica is in one of six roles:

% snippet: cpp/common/include/paglets/protocols/zombie/zombie.hpp#roles
\begin{lstlisting}[style=paglets-cpp]
enum class Role : std::uint8_t {
    recovering,  // restored after a restart; listens before it may run for election
    sleeping,    // follows a main, holds its latest checkpoint
    candidate,   // the main went silent (or a worse one appeared): asks everybody
    main,        // works; sends beats and checkpoints
    delivering,  // the work is done; delivers the result to the sink
    ended,       // the group ended here (delivered or cancelled)
};
\end{lstlisting}
% /snippet

\noindent A replica that came back from a memory image after its host
restarted is \emph{recovering}: it listens for beats first and does not
resume as main on its own. A \emph{delivering} replica is a main whose work
is done; it sends the result to the sink and resends it with every beat.

Which of two replicas is further along is decided by a \emph{rank}, a
total order on checkpoints (Listing~\ref{lst:zombie-rank}).
\code{progress} is a number the \paglet{} defines; it only grows within one
lineage, that is, along a chain of mains that each took over from the
previous one's checkpoint. \code{term} is the \emph{reign} of the main that
made the checkpoint, an epoch number in the sense of Raft's
terms~\cite{ongaro2014raft}. A replica that becomes main takes a term above every
term it has seen; its checkpoint keeps its content, gets the new reign, and
is made durable before the first beat (\code{become\_main} in
\code{zombie\_impl.hpp}). So the same content under a newer main ranks
higher, and replicas follow the newer main. Replicas that hold the same
checkpoint have equal ranks, because the author is part of the rank.

% snippet: cpp/common/include/paglets/protocols/zombie/zombie.hpp#rank [float,caption={The rank of a checkpoint: a total order on how far a lineage got},label={lst:zombie-rank}]
\begin{lstlisting}[style=paglets-cpp,float,caption={The rank of a checkpoint: a total order on how far a lineage got \src{cpp/common/include/paglets/protocols/zombie/zombie.hpp}},label={lst:zombie-rank}]
// How far a lineage got. Ranks are compared lexicographically: a finished
// lineage beats any other, then more progress, then a later reign and
// checkpoint, then the author with the lower host key ID (lower replica
// ID). The author is the main that made the checkpoint, so replicas
// holding the same checkpoint have equal ranks.
struct Rank {
    bool done = false;
    std::int64_t progress = 0;
    std::uint64_t term = 0;  // the reign of the author that made the checkpoint
    std::uint64_t seq = 0;   // checkpoint number within that reign
    ReplicaId author = 0;

    friend std::strong_ordering operator<=>(const Rank& a, const Rank& b) {
        if (auto c = a.done <=> b.done; c != 0) return c;
        if (auto c = a.progress <=> b.progress; c != 0) return c;
        if (auto c = a.term <=> b.term; c != 0) return c;
        if (auto c = a.seq <=> b.seq; c != 0) return c;
        return b.author <=> a.author;  // lower author ranks higher
    }
    friend bool operator==(const Rank&, const Rank&) = default;
};
\end{lstlisting}
% /snippet

Every role change the protocol makes is one of a fixed table of
\emph{(from, cause, to)} triples (Listing~\ref{lst:zombie-transitions},
Figure~\ref{fig:zombie-roles}). The transition function makes no other
change, and the model checker confirms that every entry occurs in some
reachable state: an unlisted change is a safety violation, and a listed one
that never occurs fails the full suite (\code{protocol\_check.cpp}).

% snippet: cpp/common/include/paglets/protocols/zombie/zombie.hpp#transitions [float,caption={Every role change of the protocol},label={lst:zombie-transitions}]
\begin{lstlisting}[style=paglets-cpp,float,caption={Every role change of the protocol \src{cpp/common/include/paglets/protocols/zombie/zombie.hpp}},label={lst:zombie-transitions}]
// Every role change the protocol makes (from, cause, to). step() makes no
// other; the model checker confirms that each one occurs.
struct Allowed {
    Role from;
    Cause cause;
    Role to;
};
inline constexpr std::array allowed_transitions = {
    Allowed{Role::sleeping, Cause::created, Role::main},
    Allowed{Role::recovering, Cause::suspect, Role::candidate},
    Allowed{Role::sleeping, Cause::suspect, Role::candidate},
    Allowed{Role::recovering, Cause::beat, Role::sleeping},
    Allowed{Role::candidate, Cause::beat, Role::sleeping},
    Allowed{Role::recovering, Cause::adopted, Role::sleeping},
    Allowed{Role::candidate, Cause::adopted, Role::sleeping},
    Allowed{Role::recovering, Cause::worse_main, Role::candidate},
    Allowed{Role::sleeping, Cause::worse_main, Role::candidate},
    Allowed{Role::main, Cause::better_main, Role::sleeping},
    Allowed{Role::candidate, Cause::round_won, Role::main},
    Allowed{Role::candidate, Cause::round_won, Role::delivering},
    Allowed{Role::candidate, Cause::round_lost, Role::sleeping},
    Allowed{Role::recovering, Cause::takeover, Role::main},
    Allowed{Role::recovering, Cause::takeover, Role::delivering},
    Allowed{Role::sleeping, Cause::takeover, Role::main},
    Allowed{Role::sleeping, Cause::takeover, Role::delivering},
    Allowed{Role::candidate, Cause::takeover, Role::main},
    Allowed{Role::candidate, Cause::takeover, Role::delivering},
    Allowed{Role::main, Cause::done, Role::delivering},
    Allowed{Role::sleeping, Cause::recovered, Role::recovering},
    Allowed{Role::recovering, Cause::recovered, Role::recovering},
    Allowed{Role::candidate, Cause::recovered, Role::recovering},
    Allowed{Role::main, Cause::recovered, Role::recovering},
    Allowed{Role::delivering, Cause::recovered, Role::recovering},
    Allowed{Role::recovering, Cause::delivered, Role::ended},
    Allowed{Role::sleeping, Cause::delivered, Role::ended},
    Allowed{Role::candidate, Cause::delivered, Role::ended},
    Allowed{Role::delivering, Cause::delivered, Role::ended},
    Allowed{Role::recovering, Cause::tomb, Role::ended},
    Allowed{Role::sleeping, Cause::tomb, Role::ended},
    Allowed{Role::candidate, Cause::tomb, Role::ended},
    Allowed{Role::main, Cause::tomb, Role::ended},
    Allowed{Role::delivering, Cause::tomb, Role::ended},
};
\end{lstlisting}
% /snippet

\begin{figure}[htbp]
\centering
\resizebox{\linewidth}{!}{%
\begin{tikzpicture}
  \node[plain] (anyrec) at (-5.0,2.4) {any role but ended};
  \node[state] (rec) at (0,2.4) {recovering};
  \node[state] (cand) at (4.4,4.2) {candidate};
  \node[state] (sleep) at (4.4,0.6) {sleeping};
  \node[state] (deliv) at (10.0,3.4) {delivering};
  \node[state] (main) at (10.0,0.6) {main};
  \node[draw=pagletsMuted,dashed,rounded corners=4pt,fit=(main)(deliv),inner sep=8pt,
        label={[event]below:leading}] (lead) {};
  \node[plain] (anyend) at (3.6,-2.2) {any role but ended};
  \node[state] (ended) at (10.0,-2.2) {ended};

  \draw[dashedarrow] (anyrec) -- node[event,above] {recovered} (rec);
  \draw[arrow] (rec) -- node[event,above,sloped,pos=0.45] {suspect, worse\_main} (cand);
  \draw[arrow] (rec) -- node[event,below,sloped,pos=0.45] {beat, adopted} (sleep);
  \draw[arrow] (sleep) to[bend right=22] node[event,right] {suspect, worse\_main} (cand);
  \draw[arrow] (cand) to[bend right=22] node[event,left,align=center] {beat, adopted,\\round\_lost} (sleep);
  \draw[arrow] (cand) -- node[event,above,sloped] {round\_won, takeover} (lead.west |- deliv);
  \draw[arrow] (sleep) -- node[event,above] {created, takeover} (lead.west |- main);
  \draw[arrow] (main.south west) to[bend left=28] node[event,below] {better\_main} (sleep.south east);
  \draw[arrow] (rec.north) to[out=70,in=110,looseness=1.1] node[event,above] {takeover} (lead.north);
  \draw[arrow] (main) -- node[event,right] {done} (deliv);
  \draw[dashedarrow] (anyend) -- node[event,above] {tomb} node[event,below] {delivered (every role but main)} (ended);
\end{tikzpicture}}
\caption{The roles and the transitions of Listing~\ref{lst:zombie-transitions}. An arrow carries the causes of the change. An arrow into the dashed box ends in \code{main}, or in \code{delivering} when the checkpoint taken over is already final. \code{recovered} leads from every role but \code{ended} to \code{recovering}; \code{tomb} leads from every role but \code{ended} to \code{ended}, \code{delivered} from every role but \code{main} and \code{ended}.}
\label{fig:zombie-roles}
\end{figure}

\subsection{Beats, elections and takeovers}\label{sec:zombies-takeover}

Replicas exchange seven messages (Table~\ref{tab:zombie-messages}). Every
message carries a term, and a receiver raises its own highest term to it.
Four timers drive the rest: \code{beat}, \code{checkpoint}, \code{suspect}
(a sleeping or recovering replica has heard no beat for a while) and
\code{round} (an election round ends; a main's pause ends).

\begin{table}[htbp]
\centering\small
\caption{The messages of the protocol (\code{zombie.hpp}).}
\label{tab:zombie-messages}
\begin{tabular}{lp{4.3cm}p{7.6cm}}
\toprule
Message & From, to & Content \\
\midrule
\code{beat} & main to every member & rank of its checkpoint, version vector of its journal, its membership \\
\code{ack} & replica to the main whose beat it heard & its rank and version vector, when it lacks something \\
\code{checkpoint} & main to one replica & the work (if the replica is behind) and the entries it lacks \\
\code{entries} & replica to main & entries the main lacks \\
\code{status?} & candidate to every member & asks for role and rank \\
\code{status} & the answer & role, rank, highest term seen \\
\code{takeover} & candidate to the best replica of its round & take over as main \\
\bottomrule
\end{tabular}
\end{table}

\textbf{Suspicion and election.} A sleeping replica that hears no beat for
\code{suspect\_ms} (10\,s) plus a random jitter of up to \code{jitter\_ms}
(5\,s) becomes a candidate and broadcasts \code{status?}. When its round
ends after \code{round\_ms} (4\,s), it picks the best of itself and every
replica that answered: the higher rank; among equal ranks a replica that
already leads; then the lower replica ID. If it is best itself, it becomes
main with a new reign. Otherwise it sends \code{takeover} to the best
replica, unless that one leads already, and follows it (\code{decide()}).
A candidate that hears a beat of a main at least as far as itself goes
back to sleep at once. A main that is asked for its status checkpoints
first and pauses its work for one round, so the rank it reported stays
true while the round is open.

\textbf{Two mains meet.} When a partition heals, each main hears the
other's beat. The one behind steps down at once: it stops its work, becomes
sleeping, acknowledges the other's checkpoint and hands over the journal
entries the other lacks. The one ahead answers with its own beat and pauses
for a round, so its rank stays what the other compares with
(\code{meet()} in \code{zombie\_impl.hpp}). The conflict is settled in at
most two exchanges, even while both keep working. Because the order on
ranks is total, two mains that meet always agree on the same one. That any
number of mains converge to one, whatever the order of meetings, is the
goal \code{one\_main} of Section~\ref{sec:zombies-model}: checked in full
for two replicas, and for three replicas only by the timed runs.

\textbf{Restarts.} After its host restarts, a replica comes back from its
memory image as \code{recovering}, whatever it was, with a new incarnation
and the working state beyond its checkpoint dropped. It listens for the
suspect timeout before it may run for election. A former main does not
resume on its own; if nobody is further along, it wins the election with
its checkpoint. After a shutdown of the whole mesh, the group is down until
hosts come back; then the most advanced replica among them continues.

\textbf{Membership growth.} The membership is a set of replica IDs that
only grows. Beats carry it, and a message from a replica not in the set
adds the sender. Every \code{expand\_ms} (30\,s) the acting main asks
\code{mesh-info} for the hosts of the mesh, clones itself to every host
that offers the \code{zombies} service and holds no replica of the group,
and adds the new replicas (input \code{AddMembers}); a host that joins the
mesh later is occupied this way within about a minute. If the mains of two
sides of a partition both make a replica on the same new host, that host's
registry refuses the second one, which ends. A host that leaves stays a
member; messages to it are lost.

\subsection{Journals and checkpoints}\label{sec:zombies-journals}

The state of a zombie has two parts. The \emph{work} is the job's state:
the winning lineage keeps its work, the others drop theirs. The
\emph{journal} holds the observations the \paglet{} made with
\code{observe()}; it is merged, never dropped. An entry is keyed by the
observing replica, its incarnation (raised on every restart, so numbers are
never reused) and a number counting from~1. Journals merge by union, which
is commutative, associative and idempotent: replicas that exchanged their
entries hold the same journal, whatever the order of exchanges. A
\emph{version vector}~\cite{parker1983versionVectors} summarizes a journal
by the gap-free prefix it holds
per observer and incarnation. A beat carries the main's version vector; a
replica that hears it asks for what it lacks (\code{ack}, answered with a
\code{checkpoint} message) and hands over what the main lacks
(\code{entries}).

A checkpoint is the main's working state and its journal, recorded with a
new rank. The main makes one, if anything changed since the last, when
the \code{checkpoint} timer fires (\code{checkpoint\_ms}, 10\,s), when it
is asked for its status and when it meets another main; and always when
the work is done. A
sleeping replica that receives a checkpoint message further along than its
own replaces its work with it; the entries are merged in any case. One rule
binds every environment that runs the state machine: \emph{durable before
visible}. The outputs of a step are carried out in order, and a
\code{Persist} output completes, the replica's state being durable, before
any later output of the same step takes effect. A main's checkpoint is
therefore stored before any replica hears of it, and a replica that adopts
a checkpoint or merges entries stores them before it answers. Not
everything a beat says is durable yet: its version vector names the
observations the main made since its last checkpoint, and its membership
the hosts it occupied since; both become durable with the main's next
checkpoint, or on a replica that asks for the entries.

\subsection{Exactly one result: the sink}\label{sec:zombies-sink}

The sink of a group is \code{host:NAME}, a host of the mesh by its ledger
name or its key ID (or a prefix of it of at least 8 hex digits). Its
\code{zombies} system \paglet{} (Section~\ref{sec:mesh}) accepts the first
result per group and sink number, stores it on that host before it
answers, and answers every later delivery with \code{accepted: false}.
Results are limited to 4\,MB: a larger one is refused by the delivering
host's own \code{zombies} service with the error \code{quota}, which the
SDK treats like an unreachable sink, so the main tries again with every
beat and the group does not end until the owner cancels it (a limit of the
current implementation). Whatever the sink answers, the group is done: the
delivering main outputs a \code{Tombstone} and ends with cause
\code{delivered}. Its host records the tombstone, and every host sends
what it knows about ended and redirected groups to every other host every
5\,s, and a new tombstone or redirect at once, so the tombstone reaches every host,
also one that was offline and comes back later; a replica learns of it from
its own host (input \code{RegistryUpdate}) and ends with cause \code{tomb}.
Tombstones are kept for 30 days. During a partition both sides may finish
and deliver; the sink accepts one, the other main learns \emph{already
delivered} and ends too.

The owner can end a group at any time (\code{paglets-host remote zombies
cancel GROUP}), which writes a tombstone of its own. If the sink cannot
come back, the owner can \emph{redirect} the group to another host, which
raises the group's \emph{sink number}. Sink numbers spread like tombstones,
every replica keeps the highest it has seen, and a delivering main delivers
to the highest it knows. A redirect is the owner's statement that the old
sink accepted nothing; if it had, two sinks hold a result. The checked
property is therefore \emph{at most one result, plus one for every
redirect}.

\subsection{The state machine and the SDK}\label{sec:zombies-sdk}

The specification is the transition function \code{Outputs step(State\&,
const Input\&)}: it applies one input to a replica and returns what the
environment must do, in order. An \code{Input} is a timer that fired, a
message that arrived,
the \paglet{}'s work (\code{Worked}, \code{Observed}), the sink's answer, a
change of the host's registry, new members, the start of the group, or a
restart (\code{Recovered}). An \code{Output} is a message to send,
\code{Persist}, a timer to set or stop, \code{StartWork} and
\code{StopWork}, \code{Deliver}, \code{Tombstone}, \code{End}, or a
\code{Transition} record. \code{State} and the messages have MessagePack
encodings.

The guest SDK runs this function in the class template
\code{paglets::patterns::Zombie<Work>}, header \code{patterns/zombie.hpp}
of the SDK. A \paglet{} derives from it and provides three things:
\code{step(work, finished)}, one step of the work on the acting main, which
calls \code{finished(done)} once, also after waiting for replies;
\code{progress(work)}, a number that only grows; and \code{result(work,
journal)}, the bytes for the sink. It may override \code{result\_name()}
(default \code{result.bin}), \code{on\_role()}, which reports every role
change, and \code{on\_zombie\_failed()}. \code{Work} is any type with \code{paglets\_encode} and \code{paglets\_decode}, for
example one from a schema header. The creator calls
\code{start\_group(sink, work, timing)}, typically from a request handler;
it starts as the main, asks \code{mesh-info} for the hosts, which must all
offer \code{zombies}, and clones itself to every other host. The Mesh
Chronicle's step asks \code{mesh-info} for the landscape and records what
changed (Listing~\ref{lst:chronicle-step}); \code{note()}, not shown,
calls \code{observe\_value()}, which puts an entry into the journal on the
acting main only. \code{finished(false)} on a failed request means: not
done, try again after \code{step\_ms}.

% snippet: cpp/examples/chronicle/chronicle.cpp#step [float,caption={The Mesh Chronicle's step: one sample of the mesh, taken on the acting main},label={lst:chronicle-step}]
\begin{lstlisting}[style=paglets-cpp,float,caption={The Mesh Chronicle's step: one sample of the mesh, taken on the acting main \src{cpp/examples/chronicle/chronicle.cpp}},label={lst:chronicle-step}]
void step(Work& w, std::function<void(bool done)> finished) override {
    auto mi = service("mesh-info");
    if (!mi) return finished(false);
    services::mesh_info::Client client{*mi};
    services::mesh_info::LandscapeRequest q;
    q.max_age_ms = std::max<std::int64_t>(3 * w.sample_ms, 3000);
    auto sent = client.landscape(q, [this, &w, finished](Result<services::mesh_info::Landscape> r, Message&) {
        if (!r) return finished(false);
        sample(w, r->hosts);
        const bool done = finish_ || now_ms() - w.started_ms >= w.lifetime_ms;
        if (done) note(Event{now_ms(), "end", here_, {}, std::to_string(w.samples) + " samples"});
        finished(done);
    });
    if (!sent) finished(false);
}
\end{lstlisting}
% /snippet

A step may wait, and the main may stop meanwhile: it met a main further
along, or a candidate's round paused it. The base class fences the step
with a generation, so a late \code{finished} of a stopped main is ignored:

% snippet: cpp/sdk/include/paglets/patterns/zombie.hpp#run_step
\begin{lstlisting}[style=paglets-cpp]
void run_step(std::uint64_t gen) {
    step(work_, [this, gen](bool done) {
        if (gen != work_gen_ || !acting()) return;  // stopped meanwhile: the result is void
        run(zp::Worked{progress(work_), done, encode(work_)});
        if (!done && gen == work_gen_ && acting()) schedule_step();
    });
}
\end{lstlisting}
% /snippet

\noindent The base class carries out the outputs. Messages go to the other
replicas through endpoints from the host's \code{zombies} registry (every
host tells the others which replica of a group it holds; apart from the
endpoint to its creator that a clone receives, no capability is passed
between replicas, and after a restart the endpoints are asked from the
registry again); timers are \code{after()} messages
to itself, fenced by a generation too; a \code{Transition} is recorded as a
trace on the host and reported through \code{on\_role()}. \code{Persist}
is where ABI v1.2 (Section~\ref{sec:lifecycle}, Appendix~\ref{app:abi})
comes in: \code{checkpoint()} stores the \paglet{}'s memory image right
after the current handler, before its next message. The base class calls it
for a \code{Persist} output and defers the remaining outputs of that step to
its next handler, which it reaches by an asynchronous message to itself
(Listing~\ref{lst:zombie-carry-out}). This is how \emph{durable before
visible} is kept in the guest.

% snippet: cpp/sdk/include/paglets/patterns/zombie.hpp#carry_out [float,caption={Durable before visible in the SDK: after a \code{Persist}, the rest of a step waits for the next handler, when the image is stored},label={lst:zombie-carry-out}]
\begin{lstlisting}[style=paglets-cpp,float,caption={Durable before visible in the SDK: after a \code{Persist}, the rest of a step waits for the next handler, when the image is stored \src{cpp/sdk/include/paglets/patterns/zombie.hpp}},label={lst:zombie-carry-out}]
// Carries out the outputs from `i` on. After a Persist the rest waits
// for the next handler, when the image is stored.
void carry_out(const zp::Outputs& outs, std::size_t i) {
    for (; i < outs.size(); ++i) {
        const auto& o = outs[i];
        if (std::holds_alternative<zp::Persist>(o)) {
            (void)checkpoint();
            deferred_.insert(deferred_.end(), outs.begin() + static_cast<std::ptrdiff_t>(i) + 1, outs.end());
            if (!flush_pending_) {
                flush_pending_ = true;
                SendOptions options;
                options.async = true;
                (void)self().send("zombie.flush", std::move(options));
            }
            return;
        }
        if (!deferred_.empty()) {  // keep the order behind earlier deferred outputs
            deferred_.push_back(o);
            continue;
        }
        apply(o);
    }
}
\end{lstlisting}
% /snippet

The base class calls \code{set\_restart(true, true)}
(Section~\ref{sec:lifecycle}) when the group starts and when a clone
arrives, so that a replica comes back from its last image after a host
restart and receives \code{on\_recovered} first. The base class answers it
with the \code{Recovered} input, drops the deferred outputs and
timers of before the restart, asks the registry for endpoints and for what
it learnt meanwhile, and the replica comes back as \code{recovering}.

\subsection{Checking the protocol}\label{sec:zombies-check}

The protocol is checked in three ways: two run the transition function,
the third checks live traces against its transition table and rank order.
The first two are run by \code{paglets-protocol-check}
(\path{cpp/tools/protocol_check}); the third runs on live meshes.

\subsubsection{The model checker}\label{sec:zombies-model}

\code{paglets::verify} (\code{verify/explorer.hpp} next to the protocol)
is an explicit-state model checker for protocols written as pure transition
functions. A model gives it a finite world: the initial state, the
successors of a state, each labelled and marked as a fault or not, a
canonical encoding, invariants and goals. \code{explore()} visits every
reachable state breadth-first, one level at a time on all hardware threads,
and checks the invariants on every state; a violation comes with the
shortest sequence of steps that leads to it. Goals are checked on the graph
of fault-free steps by backward reachability from the goal states: from
every reachable state, a goal state must be reachable without further
faults. In the terms of temporal logic this is $\mathrm{AG}\,\mathrm{EF}\,
\mathit{goal}$, with AG over every reachable state and EF over fault-free
steps only; for a single goal state it is the home-state property of
Petri-net analysis~\cite{murata1989petri}, for a set of goal states its
generalization to a home space: no dead end and no trap, so that whatever the faults did, once
they stop, the goal stays reachable, and any scheduler that takes every
enabled step with positive probability reaches it with probability~1. It is
not a liveness proof under fairness; that is future work (below).

The zombie model (\code{protocols/zombie/model.hpp}) is a few hosts, each
with one replica, its durable state and its registry entry (tombstone, sink
number); the sinks; the links; and the messages in flight, in one ordered
channel per sender and receiver, as the transport keeps one queue per peer.
Faults are steps like any other, limited by a budget per world: a host
crashes (losing what was not durable, and the messages to it); a link is
cut; a message is lost or delivered twice; a sink goes down; the owner
cancels or redirects; a timeout fires early, that is, while what it detects
is not true. Timers may fire at any time otherwise, so safety does not
depend on timing. Restarting, healing, gossip, delivery and the work are
not faults. The job in these worlds is one progress step and one
observation. Four invariants and two goals are checked
(Table~\ref{tab:zombie-properties}); Listing~\ref{lst:zombie-property}
shows one invariant.

\begin{table}[htbp]
\centering\small
\caption{The properties of the zombie model (\code{model.hpp}).}
\label{tab:zombie-properties}
\begin{tabular}{llp{8.6cm}}
\toprule
Property & Kind & Meaning \\
\midrule
\code{one\_result} & invariant & at most one sink holds a result, one more for every redirect \\
\code{best\_progress\_kept} & invariant & while the group lives, some replica durably holds the best checkpoint that was ever durable \\
\code{observations\_kept} & invariant & while the group lives, every observation that was ever durable is durable on some replica \\
\code{legal\_transitions} & invariant & every role change is in the table \\
\code{one\_main} & goal & a single main, followed by every other replica with its checkpoint, journal and membership (or the group ends) \\
\code{group\_ends} & goal & every replica ended, after a sink accepted the result (or the owner cancelled) \\
\bottomrule
\end{tabular}
\end{table}

% snippet: cpp/common/include/paglets/protocols/zombie/model.hpp#invariant [float,caption={An invariant of the model: the best durable checkpoint is never lost. \code{w.best} is a ghost variable, the best rank that was ever durable on any host},label={lst:zombie-property}]
\begin{lstlisting}[style=paglets-cpp,float,caption={An invariant of the model: the best durable checkpoint is never lost. \code{w.best} is a ghost variable, the best rank that was ever durable on any host \src{cpp/common/include/paglets/protocols/zombie/model.hpp}},label={lst:zombie-property}]
{"best_progress_kept",
 "while the group lives, some replica durably holds the best checkpoint ever made durable",
 [](const State& w) {
     if (tombed(w)) return true;
     for (const auto& h : w.hosts) {
         if (h.durable.role != Role::ended && h.durable.rank >= w.best) return true;
     }
     return false;
 }},
\end{lstlisting}
% /snippet

The canonical encoding keeps the state space finite. Terms are only
compared, raised to a maximum and incremented, so the encoding renames
them: the terms that occur anywhere in a world are mapped in order onto
small numbers, with gaps shrunk to~2 (\code{canonicalize()} in
\code{model.hpp}), and two worlds that differ only in the absolute values
of their terms get the same encoding. Messages to a replica that ended are
dropped from the encoding too, since they can have no effect. Hash
compaction keeps the visited set small: states are identified by a 128-bit
hash of their encoding~\cite{stern1995hashCompaction}; the set of visited
states holds hashes, and only the states of the current level are kept as
encodings. The probability that two distinct states of a
run share a hash is about $n^2/2^{129}$ for $n$ states, below $10^{-24}$
for the largest world below; an exact mode is planned (below).

Table~\ref{tab:zombie-results} shows what was checked. Each two-replica
world is explored completely: every reachable state within its fault
budget, with all six properties. The three-replica worlds are explored to a
depth of 11 steps from the initial state (with at most two messages in
flight per channel): the invariants hold on every state within that depth,
and the goals are not checked there, because backward reachability needs
the complete graph. Across the full suite, every entry of the transition
table occurs.

% hand-copied from docs/results/zombie_protocol.json (recorded 2026-10-13):
% the `% results: zombie_protocol` marker renders one float holding a wide
% table and thirteen per-test tables, which is too large for a page; see the
% refinement plan. Replace this table with the marker once paglets-docs-sync
% has a compact layout for LaTeX results tables. Until then keep these
% numbers in step with the JSON (checked 2026-10-14).
\begin{table}[htbp]
\centering\small
\caption{Model checking of the zombie protocol: the exhaustive worlds of the full suite, as recorded on 2026-10-13 (Darwin arm64, Release build). Every world has two replicas unless stated, a job of one step and one observation, and at most three messages in flight per channel (three replicas: two). Every named fault may occur once (two crashes: twice). Depth is the longest shortest path from the initial state, or the bound of a bounded exploration. The suite is \code{paglets-protocol-check --suite full}; the numbers are those of \code{docs/results/zombie\_protocol.json}.}
\label{tab:zombie-results}
\begin{tabular}{lrrrr}
\toprule
World & States & Steps & Depth & Time (s) \\
\midrule
crash, cut & 12\,908\,826 & 50\,102\,693 & 47 & 565.9 \\
crash, loss & 3\,961\,354 & 17\,259\,715 & 43 & 182.8 \\
crash, duplicate & 4\,380\,410 & 18\,416\,904 & 40 & 205.8 \\
crash, early timeout & 7\,807\,220 & 30\,116\,664 & 47 & 308.0 \\
two crashes & 22\,356\,397 & 81\,063\,681 & 53 & 519.2 \\
cut, loss, early timeout & 4\,000\,767 & 19\,983\,567 & 41 & 85.1 \\
crash, cancel & 3\,935\,759 & 19\,606\,519 & 41 & 54.9 \\
sink down, redirect & 378\,252 & 1\,698\,056 & 30 & 5.0 \\
one joining later: crash, cut & 13\,205\,168 & 51\,296\,859 & 49 & 215.5 \\
one joining later: crash, early timeout & 8\,051\,595 & 31\,115\,057 & 49 & 102.6 \\
\midrule
three replicas: crash, cut (bounded) & 9\,188\,336 & 24\,420\,885 & 11 & 79.3 \\
three replicas: two crashes (bounded) & 5\,461\,027 & 14\,745\,614 & 11 & 40.7 \\
\bottomrule
\end{tabular}
\end{table}

\subsubsection{Timed runs}\label{sec:zombies-sim}

The exhaustive check abstracts time away. A simulator
(\code{protocols/zombie/sim.hpp}) adds it back: the same \code{step()}
under virtual clocks, with the real timer durations, message latency of
1--50\,ms, a worker that makes one progress step per second for 300 steps
and an observation every ten, and random faults: crashes with a restart
after 5--60\,s, cuts healed after 5--90\,s, 2\% loss and 1\% duplication,
during a fault phase of 180\,s with a mean of 15\,s between faults. Three
replicas take part, one on a host that joins during the run. Each run has
a sink that, like the real one, accepts the first result and refuses later
ones; the run checks that a result was delivered, that every replica ends,
that the delivered journal holds every observation that was durable anywhere at
delivery, that every role change is in the table, and that a takeover
after the main crashed, with nothing else happening until a main worked
again, took at most $\mathit{suspect} + \mathit{jitter} + 2\,\mathit{round}
+ \mathit{beat} + 4\,\mathit{latency}_{\max}$, which is 25.2\,s with the
defaults. Every run is reproducible from its seed (\code{--trace SEED}
prints its role changes and faults). The 10\,000 runs recorded with
Table~\ref{tab:zombie-results} contain 119\,289 faults and 3\,370 measured
takeovers, the longest of 18.9\,s, and no failure. These runs sample
instead of enumerating, and they check a different, partly weaker, set of
properties than the model; they are an addition, not the proof.

\subsubsection{Live conformance}\label{sec:zombies-traces}

Replicas record their protocol steps on their hosts through the
\code{zombies} service: every role change with its cause and the replica's
rank, and every answer of a sink. \code{paglets-host remote zombies check
GROUP} collects these records from every live host and checks them against
the specification (\code{protocols/zombie/trace.hpp}): every role change is
in the table, no replica's rank falls except across a restart, nothing is
recorded after a replica ended, and at most one result was accepted, one
more for every redirect. This check connects the live system to the model,
and it marks what the model leaves out. The model carries the protocol's
messages but not the transport that moves them; it reduces the hosts'
gossip to a merge of tombstone and sink number between two live hosts, the
\code{zombies} service to a registry entry per host and an atomic register
per sink, and the SDK's carrying out of outputs (timers, endpoints, clones
to new hosts, the memory image) to steps of its own. The trace check, in
turn, sees only role changes, ranks and sink answers, not the messages. The
in-process tests and the live run of the Mesh Chronicle
(Section~\ref{sec:evaluation}) end with this check, and passed it.

\subsubsection{What is verified, and what is not}\label{sec:zombies-claims}

To state it precisely: in every two-replica world of
Table~\ref{tab:zombie-results}, the four invariants of
Table~\ref{tab:zombie-properties} hold in every reachable state, and from
every reachable state each of the two goals is reachable over fault-free
steps, under the fault budgets of those worlds (for example one crash and one cut, or two
crashes), a job of one step and one observation, at most three messages in
flight per channel, and up to the hash compaction's omission probability.
The four invariants hold in every state within 11 steps of the three-replica
worlds. Longer jobs, more faults and three or more replicas over their
whole state space are covered only by the timed runs and the live checks.
The model has limits of its own: it reduces the registry's gossip to a
merge step between two live hosts, and it treats a sink as a single atomic
register.

Several tightenings are planned so that these claims become sharp. An
exact mode of the checker without hash compaction for the two-replica
worlds, so that those results are exact and the omission probability is
quoted only for larger worlds. A fair-cycle check (every bottom strongly
connected component of the fault-free graph contains a goal state), which
gives liveness under weak fairness next to the recoverability result. A
mutation suite: known bugs seeded into the protocol (a dropped hand-over,
a stale beat accepted, the sink number skipped on a redirect), each shown
to be caught with a counterexample. A written soundness argument for the
term renaming, with an assertion in the model that a transition never does
more with a term than compare, maximize and increment. And three replicas
in full, by partial-order reduction or disk-backed state storage, or a
cutoff argument that $n$ replicas behave like two or three.
